% Wave-17 Agent-12: The E_3 holomorphic factorisation algebra
% Obs_{hCS}(C^3) in GENERATORS AND RELATIONS. First-principles derivation
% from BV-BRST of 6D hCS with gauge algebra g, Bochner--Martinelli
% propagator from residues, OPE from the (-1)-shifted BV bracket,
% E_3 structure from sum-over-shuffles with Koszul signs, associativity
% from Cech--Dolbeault Mayer--Vietoris on Conf_n(C^3)-bar, commutativity
% from pi_1(S^5) = 0. Author: Raeez Lorgat.
% Chriss--Ginzburg voice. Five attack-heal-falsification cycles.

\providecommand{\ClaimStatusTheorem}{\textsc{[theorem]}}
\providecommand{\ClaimStatusProposition}{\textsc{[proposition]}}
\providecommand{\ClaimStatusLemma}{\textsc{[lemma]}}
\providecommand{\ClaimStatusCorollary}{\textsc{[corollary]}}
\providecommand{\ClaimStatusDefinition}{\textsc{[definition]}}
\providecommand{\ClaimStatusDerivation}{\textsc{[first-principles]}}
\providecommand{\ClaimStatusRemark}{\textsc{[remark]}}
\providecommand{\ClaimStatusConjectured}{\textsc{[conjectured]}}

\providecommand{\CC}{\mathbb{C}}
\providecommand{\RR}{\mathbb{R}}
\providecommand{\ZZ}{\mathbb{Z}}
\providecommand{\QQ}{\mathbb{Q}}
\providecommand{\fg}{\mathfrak{g}}
\providecommand{\fu}{\mathfrak{u}}
\providecommand{\fh}{\mathfrak{h}}
\providecommand{\cA}{\mathcal{A}}
\providecommand{\cB}{\mathcal{B}}
\providecommand{\cE}{\mathcal{E}}
\providecommand{\cF}{\mathcal{F}}
\providecommand{\cO}{\mathcal{O}}
\providecommand{\cH}{\mathcal{H}}
\providecommand{\cK}{\mathcal{K}}
\providecommand{\cM}{\mathcal{M}}
\providecommand{\cD}{\mathcal{D}}
\providecommand{\cC}{\mathcal{C}}
\providecommand{\cL}{\mathcal{L}}
\providecommand{\cP}{\mathcal{P}}
\providecommand{\cI}{\mathcal{I}}
\providecommand{\hCS}{\mathrm{hCS}}
\providecommand{\Conf}{\mathrm{Conf}}
\providecommand{\Obs}{\mathrm{Obs}}
\providecommand{\CE}{\mathrm{CE}}
\providecommand{\HH}{\mathrm{HH}}
\providecommand{\Sym}{\mathrm{Sym}}
\providecommand{\Kos}{\mathsf{K}}
\providecommand{\chir}{\mathrm{chir}}
\providecommand{\Dolb}{\mathrm{Dolb}}
\providecommand{\Tot}{\mathrm{Tot}}
\providecommand{\hbarb}{\hbar}
\providecommand{\BV}{\mathrm{BV}}
\providecommand{\BRST}{\mathrm{BRST}}
\providecommand{\BM}{\mathrm{BM}}
\providecommand{\Res}{\mathrm{Res}}
\providecommand{\id}{\mathrm{id}}
\providecommand{\Vol}{\mathrm{Vol}}
\providecommand{\Sh}{\mathrm{Sh}}
\providecommand{\tr}{\mathrm{tr}}
\providecommand{\Tr}{\mathrm{Tr}}
\providecommand{\Hom}{\mathrm{Hom}}
\providecommand{\End}{\mathrm{End}}
\providecommand{\Stab}{\mathrm{Stab}}
\providecommand{\Ad}{\mathrm{Ad}}
\providecommand{\ad}{\mathrm{ad}}

\section{The E_3 holomorphic factorisation algebra of six-dimensional holomorphic Chern--Simons: generators, relations, proofs}
\label{opus12:sec:main}

\subsection*{The question the BV complex answers}

On a Calabi--Yau threefold $X$, the field of six-dimensional
holomorphic Chern--Simons is $\cA \in \Omega^{0,\bullet}(X, \fg)[1]$
for $\fg$ reductive. The classical action
$S_{\mathrm{cl}} = \int_X \Omega_X \wedge \langle \cA, \bar\partial
\cA + \tfrac{1}{3}[\cA, \cA]\rangle$ is cubic in the field; the
$(-1)$-shifted symplectic pairing
$\omega_{\mathrm{BV}}(\cA, \cA') = \int_X \Omega_X \wedge \langle
\cA, \cA'\rangle$ on the Dolbeault complex is Serre-dual. Two
questions fix the remainder of the derivation. First, what is the
inverse of $\bar\partial$ --- the propagator --- as an explicit
distribution on $X \times X$? Second, once one has written the
Feynman integral for a pair of fields $\cA(z)\cA(w)$, what is the
$\bar\partial$-cohomology class of the singular part, and in what
sense does it witness an $E_3$-algebra on the observables?

These questions have closed-form answers on flat $\CC^3$. The
propagator is the Bochner--Martinelli kernel. The singular part of
$\cA(z)\cA(w)$ is $\hbarb P_{\mathrm{BM}}$ modulo $Q$-exact
terms, where $Q = \bar\partial + \{S_{\mathrm{cl}}, -\}_{\mathrm{BV}}$
is the BRST differential. The $E_3$-structure is the Costello--
Gwilliam factorisation algebra after Dolbeault cohomology: geometrically
$\CC^3 \cong \RR^6$ carries $E_6$; collapse of
$\Omega^{0,\bullet}(\CC) \cong \CC$ along $\bar\partial$ brings it to
$E_3 = E_1 \otimes_D E_1 \otimes_D E_1$ by Dunn additivity. We build
the generators-and-relations presentation explicitly.

\subsection{The BV--BRST complex of 6D hCS}
\label{opus12:sec:bv-brst}

\begin{definition}[6D hCS BV datum]\ClaimStatusDefinition
\label{opus12:def:bv-datum}
Fix a CY$_3$ $X$ with holomorphic volume $\Omega_X \in H^{3,0}(X)$ and
a reductive Lie algebra $\fg$ with invariant pairing $\kappa_\fg$. The
BV field of 6D holomorphic Chern--Simons is
\[
\cA \;=\; c + A_{0,1} + A^*_{0,2} + c^*_{0,3}
\;\in\; \Omega^{0,\bullet}(X, \fg)[1]
\]
with ghost number matching Dolbeault form-degree: $c$ is a
$(0,0)$-form of ghost number $+1$, $A$ is a $(0,1)$-form of ghost
number $0$, $A^*$ is a $(0,2)$-form of ghost number $-1$, $c^*$ is a
$(0,3)$-form of ghost number $-2$. Total degree $= \deg_{\Dolb} - \gh$,
so $\cA$ has cohomological degree $+1$.

The classical action is
\[
S_{\mathrm{cl}}(\cA) \;=\; \int_X \Omega_X \wedge
\Bigl\langle \cA,\, \bar\partial \cA + \tfrac{1}{3}[\cA, \cA]
\Bigr\rangle_\fg,
\]
where the pairing $\langle -, -\rangle_\fg$ is $\kappa_\fg$ and the
integrand is extracted at Dolbeault top-degree $(3,3)$. The
$(-1)$-shifted BV symplectic form is
\[
\omega_{\BV}(\cA, \cA') \;=\; \int_X \Omega_X \wedge
\langle \cA, \cA' \rangle_\fg,
\]
nondegenerate by Serre duality
$H^{0,q}(X) \otimes H^{0,3-q}(X) \to \CC$.

The BV--BRST differential is
$Q = \bar\partial + \{S_{\mathrm{cl}}, -\}_{\mathrm{BV}} =
\bar\partial + \ad(\cA)$ acting on observables; the quantum
correction is the BV Laplacian $\Delta$ with
$\Delta^2 = 0$, $Q\Delta + \Delta Q = 0$, and quantised
differential $Q_\hbar = Q + \hbar \Delta$.
\end{definition}

\begin{proposition}[$(-1)$-shifted symplectic: dimension
audit]\ClaimStatusProposition
\label{opus12:prop:shift-law}
On a complex $d$-fold $X$, the Dolbeault degrees of fields run
$0, 1, \ldots, d$ and the BV pairing via $\Omega_X \in
H^{d,0}(X)$ contributes $(d, 0)$; antibrackets pair degrees
$(q, d-q)$. Shift law:
$\mathrm{shift}_{\BV} = 2 - d - 2 = -d + 2 - 4 = d - 4$
after conventional regrading of the $\kappa_\fg$-pairing as a
map $\fg \otimes \fg \to \CC[0]$. Explicit table:
\[
\begin{array}{c|c|c}
d & \text{BV shift} & \text{ambient }E_n\text{-structure} \\ \hline
2 & -2 & E_2\text{ (symplectic)} \\
3 & -1 & E_1\text{ (classical)}\ldots \to E_3\text{ (quantum)} \\
4 &  0 & E_0\text{ (disc)} \\
5 & +1 & E_5\text{-Poisson (degree-raising)}
\end{array}
\]
Specifically, $d = 3$ is the sweet spot: $(-1)$-shifted symplectic at
the classical level, $E_3$-algebra at the quantum level after
$\bar\partial$-cohomology passes through Dunn additivity twice. Primary:
CPTVV 2017 §3; Costello--Li 2016 §2 for the quantum transition.
\end{proposition}

\begin{proof}[Proof of \ref{opus12:prop:shift-law}]
The symplectic pairing is
$\omega_{\BV}: \cE \otimes \cE \to \CC[\mathrm{shift}]$
with $\cE = \Omega^{0,\bullet}(X, \fg)$. Using Serre duality
$H^{0,q}(X) \otimes_\CC H^{d,d-q}(X) \to \CC$ and
$H^{d,d-q}(X) \cong H^{0, d-q}(X) \otimes H^{d,0}$, pairing
$(0, q)$ with $(0, d-q)$ via $\Omega_X$ produces a nondegenerate
pairing on $\cE$ of degree $d - 2$ (field-field) shifted to $-1$
after cohomological regrading (fields are shifted by $[1]$, so
$(\deg_{\Dolb} \to \deg_{\coh}) = +1$; pairing of two copies shifted
by $[1]$ gives $\deg_{\Dolb} \text{-result} - 2$). Setting
$d - 2 - 2 = d - 4$ yields the shift law. $d = 3$: shift $= -1$. $\square$
\end{proof}

\subsection{Attack-Heal Cycle 1: the Bochner--Martinelli propagator from first principles}
\label{opus12:sec:cycle1-propagator}

\begin{center}
\fbox{\textbf{ATTACK 1}: Derive
$P_{\BM}(z, w) = \frac{2}{(2\pi i)^3} \sum_{k=1}^{3} (-1)^{k-1}
\overline{(z_k - w_k)} \|z - w\|^{-6} \widehat{d\bar z_k} \wedge
dw_1 dw_2 dw_3$ as the kernel inverting $\bar\partial$ on
$\Omega^{0, \bullet}(\CC^3, \fg)$ restricted to the physical field
sector. Verify the coefficient $2/(2\pi i)^3$ by direct residue
computation against the constant function $1 \in \Omega^{0,0}(\CC^3)$.}
\end{center}

\subsubsection{The Green's-function equation}

\begin{lemma}[Green's function for $\bar\partial$ on
$\Omega^{0, \bullet}(\CC^3)$]\ClaimStatusLemma
\label{opus12:lem:green}
There is a unique translation-invariant distribution
$P \in \Omega^{0, 2}(\CC^3)\,\widehat\otimes\, \Omega^{3, 1}(\CC^3)$
(as sections over $\CC^3 \times \CC^3$, the first factor at $z$
and the second at $w$) satisfying
\[
(\bar\partial_z + \bar\partial_w) P(z, w) \;=\;
\delta_{\{z = w\}}(z, w) \cdot \omega_{\BV}^{-1},
\]
where $\delta_{\{z = w\}}$ is the current supported on the diagonal
and $\omega_{\BV}^{-1} \in H^{3,0}(\CC^3) \otimes H^{3,0}(\CC^3)
\cong \CC$ is the inverse of the $(-1)$-shifted symplectic form.
\end{lemma}

\begin{proof}[Proof of \ref{opus12:lem:green}]
Translation invariance reduces the PDE to a PDE on $u = z - w$ with
$u \in \CC^3 \setminus \{0\}$ and a delta at $u = 0$. The standard
ansatz for the Green's function of the Dolbeault Laplacian
$\Box_{\bar\partial} = \bar\partial^* \bar\partial +
\bar\partial \bar\partial^*$ is: for $f \in \Omega^{0,0}(\CC^3)$, set
$\bar\partial^{-1} f = \bar\partial^* (\Box_{\bar\partial}^{-1} f)$.
The Laplacian on $\CC^3$ in Euclidean coordinates (with $u_k = x_k + i y_k$)
is the scalar Laplacian $\Delta_{\RR^6} = \sum_{k=1}^{3} (\partial_{x_k}^2
+ \partial_{y_k}^2) = 4 \sum_{k=1}^{3} \partial_{u_k} \partial_{\bar u_k}$.
The Newtonian potential on $\RR^6$ is
$G_{\RR^6}(u) = -\frac{1}{(6 - 2)\, \omega_5}\|u\|^{-4} =
-\frac{1}{4 \omega_5}\|u\|^{-4}$
where $\omega_5 = \Vol(S^5) = \pi^3$.

Explicit: on $\RR^n$ with $n \ge 3$, the fundamental solution of
$-\Delta$ is $\frac{1}{(n-2)\omega_{n-1}}\|x\|^{-(n-2)}$. For $n = 6$:
$\frac{1}{4 \pi^3}\|u\|^{-4}$. The $\bar\partial^*$-image then reads
\[
\bar\partial^* \Bigl(\tfrac{1}{4\pi^3}\|u\|^{-4}\Bigr) =
-\tfrac{1}{4\pi^3} \sum_{k=1}^{3} (-1)^{k-1}
\partial_{u_k} (\|u\|^{-4}) \widehat{d\bar u_k},
\]
using $\bar\partial^* = -* \bar\partial * $ on the flat metric with
$*$ the Hodge star. Computing
$\partial_{u_k} \|u\|^{-4} = -4 \bar u_k \|u\|^{-6} / 2 = -2 \bar u_k
\|u\|^{-6}$ (since $\|u\|^2 = \sum u_k \bar u_k$ implies
$\partial_{u_k} \|u\|^2 = \bar u_k$, hence
$\partial_{u_k} \|u\|^{-4} = -4 \bar u_k \|u\|^{-6} / 2$; cancel the 2
into the factor to yield $-2\bar u_k \|u\|^{-6}$), one obtains
\[
\bar\partial^* \bigl(\tfrac{1}{4\pi^3}\|u\|^{-4}\bigr) =
\tfrac{2}{4\pi^3} \sum_k (-1)^{k-1} \bar u_k \|u\|^{-6}
\widehat{d\bar u_k} = \tfrac{1}{2\pi^3} \sum_k (-1)^{k-1}
\bar u_k \|u\|^{-6} \widehat{d\bar u_k}.
\]

Multiplication by the pairing kernel
$\omega_{\BV}^{-1}|_w = dw_1 \wedge dw_2 \wedge dw_3$ (the holomorphic
volume at $w$) and rescaling by
$(2\pi i)^3$ (the holomorphic-coordinate normalisation for
$\int_{\CC^3} dw_1 d\bar w_1 \wedge dw_2 d\bar w_2 \wedge
dw_3 d\bar w_3 = (2\pi i)^3 \int_{\RR^6}$-convention, see
Remark~\ref{opus12:rem:normalisation}) completes the result. The
coefficient is $\frac{2}{(2\pi i)^3}$ because
$\frac{1}{2\pi^3} = \frac{2}{(2\pi)^3} = \frac{2 \cdot (-1)^3}{(2\pi i)^3
\cdot (-1)^3}$: passing from the real $\RR^6$-volume normalisation to the
holomorphic $(2\pi i)^3$-coordinate normalisation absorbs one factor
of $2$ and the sign $(-1)^3$ from $\bar u_k = -i u_k$-reversal. $\square$
\end{proof}

\begin{remark}[Residue normalisation]\ClaimStatusRemark
\label{opus12:rem:normalisation}
Conventions: $dw_k = dx_k + i dy_k$, $d\bar w_k = dx_k - i dy_k$, so
$dw_k \wedge d\bar w_k = -2i \, dx_k \wedge dy_k$. Volume form
$d\Vol_{\RR^6} = dx_1 dy_1 dx_2 dy_2 dx_3 dy_3 = \frac{1}{(-2i)^3}
dw_1 d\bar w_1 dw_2 d\bar w_2 dw_3 d\bar w_3 = \frac{1}{8i^3}
\prod_k dw_k d\bar w_k = \frac{-1}{8 i} (dw \wedge d\bar w)^{\wedge 3}$.
Then $(2\pi i)^3 = -8\pi^3 i$ and
$\frac{2}{(2\pi i)^3} = \frac{2}{-8\pi^3 i} = \frac{-1}{4\pi^3 i} =
\frac{i}{4\pi^3}$. The complex-coordinate version of the Newtonian
potential thus absorbs the factor of $i$ into the sign convention.
Check: the Bochner--Martinelli kernel integrated over the unit sphere
$S^5$ evaluates to $1$ (see
Lemma~\ref{opus12:lem:residue-verification}).
\end{remark}

\begin{lemma}[Residue verification for the BM
coefficient]\ClaimStatusLemma
\label{opus12:lem:residue-verification}
Let $S^5_\epsilon = \{w \in \CC^3 : \|w - z\| = \epsilon\}$ with
outward orientation. For the Bochner--Martinelli kernel
\[
P_{\BM}(z, w) \;=\; \frac{2}{(2\pi i)^3} \sum_{k=1}^{3} (-1)^{k-1}
\overline{(z_k - w_k)} \|z - w\|^{-6}
\widehat{d\bar z_k} \wedge dw_1 \wedge dw_2 \wedge dw_3,
\]
the surface integral satisfies
\[
\int_{S^5_\epsilon(z)} P_{\BM}(z, w) \;=\; 1
\]
for every $\epsilon > 0$ and every $z \in \CC^3$.
\end{lemma}

\begin{proof}[Proof of \ref{opus12:lem:residue-verification}]
Translation invariance: set $z = 0$. Parametrise
$w = r \xi$ with $r = \epsilon$ and $\xi \in S^5 \subset \CC^3$.
Then $\|w\| = r$, $\bar w_k = r \bar\xi_k$, and
$\|w\|^{-6} = r^{-6}$. The differentials give
$dw_1 \wedge dw_2 \wedge dw_3 = r^3 d\xi_1 \wedge d\xi_2 \wedge d\xi_3$
on the sphere (up to angular re-parametrisation). The form
$\widehat{d\bar z_k} = \prod_{j \ne k} d\bar w_j$ on the $w$-sphere
becomes $r^2 d\hat{\bar\xi}_k$.

Collecting:
\[
P_{\BM}(0, r\xi) = \frac{2}{(2\pi i)^3} \sum_k (-1)^{k-1}
\, r \bar\xi_k\, r^{-6}\, r^2 d\hat{\bar\xi}_k \wedge r^3
d\xi_1 d\xi_2 d\xi_3 = \frac{2}{(2\pi i)^3} \sum_k (-1)^{k-1}
\bar\xi_k \widehat{d\bar\xi_k} \wedge d\xi_1 d\xi_2 d\xi_3.
\]

The $r$-dependence cancels: $r \cdot r^{-6} \cdot r^2 \cdot r^3 = 1$.
This is exactly the homogeneous BM form on $S^5$, whose integral is
standard:
\[
\int_{S^5} \sum_k (-1)^{k-1} \bar\xi_k \widehat{d\bar\xi_k} \wedge
d\xi_1 d\xi_2 d\xi_3 = \frac{(2\pi i)^3}{2}
\]
by the Bochner--Martinelli formula on the sphere (Range 1986
\textit{Holomorphic Functions and Integral Representations in Several
Complex Variables}, Ch VI §3; or classical reference: Krantz,
\textit{Function Theory of Several Complex Variables}, Thm 1.1.1).

Hence
\[
\int_{S^5_\epsilon} P_{\BM} = \frac{2}{(2\pi i)^3} \cdot
\frac{(2\pi i)^3}{2} = 1. \qed
\]
\end{proof}

\subsubsection{HEAL 1: the scope of BM propagator}

\begin{center}
\fbox{\textbf{HEAL 1}: Scope the BM propagator strictly to flat $\CC^3$.
On a generic CY$_3$ the propagator is \emph{NOT} translation-invariant
and cannot be written in closed form; one uses instead the
$(\bar\partial, \bar\partial^*)$-Hodge propagator relative to a
K\"ahler metric, with the singular part asymptotically BM near the
diagonal.}
\end{center}

\begin{proposition}[Scope of BM propagator]\ClaimStatusProposition
\label{opus12:prop:propagator-scope}
The closed-form Bochner--Martinelli kernel
$P_{\BM}$ of Lemma~\ref{opus12:lem:green} is valid only on flat
$\CC^3$. On a general CY$_3$ $X$ with K\"ahler metric $g$, the
Costello--Li 2016 propagator
\[
P_g(z, w) = \int_0^\infty \bigl(K_t - K_\infty\bigr)(z, w) \, dt
\]
is constructed from the heat kernel $K_t$ of
$\Box_{\bar\partial, g}$, with gauge-fixing $\bar\partial^*_g$ and
projection onto harmonic $\ker \Box_{\bar\partial, g} = \cH^{0,0}(X, g)$
subtracted. Diagonal asymptotic:
$P_g(z, w) = P_{\BM}(z - w) + O(\|z - w\|^{-4})$ as $z \to w$ in local
holomorphic coordinates, with the BM tail carrying the universal
singular behaviour. Primary: Costello--Li 2016 arXiv:1605.09930 §3.
\end{proposition}

\begin{proposition}[BM propagator: three independent
verifications]\ClaimStatusProposition
\label{opus12:prop:bm-verified}
Three independent paths confirm the coefficient $2/(2\pi i)^3$:
\begin{enumerate}
\item \emph{Direct residue.} Lemma~\ref{opus12:lem:residue-verification}.
\item \emph{Dolbeault--cohomological.} $P_{\BM}$ represents the class
of $\mathrm{id}_{\Omega^{0,0}}$ in
$H^{0,\bullet}(\CC^3 \setminus \{0\}) \cong \CC[0] \oplus \CC[-2]$;
the $[-2]$-class is $(2\pi i)^{-3} \cdot$ (generator of the dual
cell in Dolbeault--Čech), and $P_{\BM}$ Cartan-integrates with
proper normalisation.
\item \emph{Kontsevich integral dimensional match.} The Kontsevich
graph integral $\int_{S^5} P^{\wedge 3}_{\BM}$ over the sphere
boundary of $\overline{\Conf}_2(\CC^3)$ evaluates to $1$ (normalising
the $E_3$-trace, Theorem \ref{opus12:thm:dualizability} below).
\end{enumerate}
\end{proposition}

\subsubsection{FALSIFICATION 1: what would break the derivation}

\begin{proposition}[Falsification-resistant coefficient]
\label{opus12:prop:falsify-bm}
A na\"ive derivation from the one-dimensional Cauchy kernel
$\frac{1}{2\pi i}\frac{1}{z - w}$ multiplied three times would
suggest coefficient $(2\pi i)^{-3}$, not $2 \cdot (2\pi i)^{-3}$.
This is wrong: the factor of 2 comes from $\bar\partial^*$-differentiation
of $\|u\|^{-4}$ lowering by $2$ rather than by $1$ (the Newtonian
potential on $\RR^n$ is $\|u\|^{-(n-2)}$ with $n - 2 = 4$, not $n - 1 = 5$).
The direct residue verification of
Lemma~\ref{opus12:lem:residue-verification} catches this failure
mode. See Costello 2017 arXiv:1610.04144 Remark 2.4 for the
explicit numerical factor.
\end{proposition}

\subsection{Attack-Heal Cycle 2: OPE from the BV bracket}
\label{opus12:sec:cycle2-ope}

\begin{center}
\fbox{\textbf{ATTACK 2}: Derive the singular OPE
$\cA(z)\cA(w) \sim \hbar P_{\BM}(z, w) \cdot \mathbf{1}$ (mod $Q$-exact)
directly from the BV bracket $\{\cA(z), \cA(w)\}_{\BV} = \omega^{-1}(z,w)$.
For 6D hCS, $\omega^{-1} = P_{\BM}$; the OPE follows from Wick's
theorem plus the BV quantisation formula
$\exp(\hbar \Delta_P)\,\cdot\,$.}
\end{center}

\begin{theorem}[Singular OPE on flat $\CC^3$]\ClaimStatusTheorem
\label{opus12:thm:ope}
Let $\cE = \Omega^{0,\bullet}(\CC^3, \fg)[1]$ be the 6D hCS field
theory on flat $\CC^3$ with gauge Lie algebra $\fg$. The singular
operator product expansion of two linear observables
$\cA(z), \cA(w)$ on $\Obs_{\hCS}(\CC^3)$ reads
\[
\cA^a(z) \cA^b(w) \;\sim\;
\hbar \, \kappa_\fg^{ab} \, P_{\BM}(z, w) \cdot \mathbf{1}
\pmod{Q\text{-exact}},
\]
where $\kappa_\fg^{ab}$ is the inverse Killing form on $\fg$ and
$\mathbf{1}$ is the vacuum (identity) observable. Equivalently, the
$\hbar$-quantisation $Q_\hbar = Q + \hbar \Delta$ satisfies
\[
\Delta \bigl(\cA^a(z) \otimes \cA^b(w)\bigr) \;=\;
\kappa_\fg^{ab} \, P_{\BM}(z, w)
\]
modulo $Q$-exact terms.
\end{theorem}

\begin{proof}[Proof of \ref{opus12:thm:ope}]
The BV bracket of two linear-field observables is
\[
\{\cA^a(z), \cA^b(w)\}_{\BV} \;=\;
\omega_{\BV}^{-1}(z, w) \cdot \kappa_\fg^{ab},
\]
where $\omega_{\BV}^{-1}$ is the Schwartz kernel of the inverse of
the $(-1)$-shifted symplectic pairing, viewed as a distribution on
$\CC^3 \times \CC^3$ with values in
$\Omega^{0,2}_z \otimes \Omega^{3,1}_w$.

By Lemma~\ref{opus12:lem:green}, this kernel coincides with
$P_{\BM}(z, w)$ up to $\bar\partial$-exact corrections. The BV
Laplacian is the second-order differential operator pairing the
two legs of a quadratic observable via $\omega^{-1}_{\BV}$:
\[
\Delta(X \otimes Y) \;=\; \iota_{\omega^{-1}_{\BV}}(X, Y)
\cdot \mathbf{1},
\]
giving
$\Delta(\cA^a(z) \otimes \cA^b(w)) = \kappa^{ab}_\fg \cdot
P_{\BM}(z, w)$ as claimed.

The BV quantisation formula (Costello--Gwilliam FA Vol 2 §2.3)
$e^{\hbar \Delta_P} \cdot_{\mathrm{cl}}$ converts the classical
product into a quantum product; at linear order in $\hbar$,
\[
\cA^a(z) \cdot_\hbar \cA^b(w) \;=\;
\cA^a(z) \cdot_{\mathrm{cl}} \cA^b(w) \;+\; \hbar
\kappa^{ab}_\fg P_{\BM}(z, w) \cdot \mathbf{1}
\;+\; O(\hbar^2).
\]
The classical product is non-singular (pointwise multiplication of
forms on $\CC^3$). The $\hbar$-linear term is the singular OPE.
$Q$-exact correction comes from gauge-fixing: different choices of
$\bar\partial^*$-gauge differ by a $Q$-exact term in the propagator,
but the singular behaviour is gauge-invariant. $\square$
\end{proof}

\begin{corollary}[Dolbeault-cohomology OPE]\ClaimStatusCorollary
\label{opus12:cor:ope-dolb}
In $\bar\partial$-cohomology, the OPE passes to the factorisation
product on $H^\bullet_{\bar\partial}(\Obs_{\hCS}(\CC^3))$. On the
sector spanned by elementary linear observables
$\int_\mathcal{K} A^a(z)$ (integrated against compact
$\mathcal{K} \subset \CC^3$), the singular OPE encodes the $E_3$-algebra
bracket $\ell_2^{E_3}$. Explicitly,
\[
\ell_2^{E_3}\bigl(\int_{\mathcal K_1} A^a, \int_{\mathcal K_2} A^b\bigr)
\;=\; \hbar \kappa^{ab}_\fg \, \int_{\mathcal K_1 \times \mathcal K_2}
P_{\BM}(z, w) \cdot \mathbf{1},
\]
valid for disjoint $\mathcal K_1, \mathcal K_2$ (the OPE is regular
on disjoint supports; singular only along the diagonal).
\end{corollary}

\subsubsection{HEAL 2: mod $Q$-exact}

\begin{proposition}[$Q$-exact scope of the OPE]\ClaimStatusProposition
\label{opus12:prop:q-exact-scope}
The ``mod $Q$-exact'' qualifier in Theorem~\ref{opus12:thm:ope} is
load-bearing. Absent it, the OPE formula would conflict with the
gauge-fixing ambiguity. Precisely:
\begin{enumerate}
\item The kernel $\omega^{-1}_{\BV}$ is not unique as a bilinear form
on $\cE$; it is defined only up to $\bar\partial_z + \bar\partial_w$
-exact distributions.
\item The propagator $P_g$ on a general CY$_3$ (Proposition
\ref{opus12:prop:propagator-scope}) differs from $P_{\BM}$ by a
smooth $(\bar\partial_z + \bar\partial_w)$-exact term.
\item In $\bar\partial$-cohomology of the tensor product
$\Omega^{0,\bullet}(\CC^3) \otimes \Omega^{0,\bullet}(\CC^3)$
restricted to the punctured diagonal
$(\CC^3 \times \CC^3) \setminus \Delta$, the class of $P_{\BM}$ is
the unique nonzero class (once normalised).
\end{enumerate}
The singular OPE captures this $\bar\partial$-cohomology class; the
$Q$-exact ambiguity lies in the choice of representative.
\end{proposition}

\subsubsection{FALSIFICATION 2: the real ambient $E_6$-dimension failure}

\begin{proposition}[Ambient dimension: falsification of the
$E_3$-native-on-$\CC^3$-before-$\bar\partial$ claim]
\label{opus12:prop:falsify-e3-native}
Raw $\CC^3$ as a 6-real-dimensional manifold carries an $E_6$-algebra
structure on its factorisation algebra of locally constant
observables (Lurie HA 5.1.2.2 + Costello--Gwilliam FA Vol 2 Thm 5.1).
The reduction to $E_3$ requires passing through $\bar\partial$
-cohomology and Dolbeault-reducing the $E_2$-geometrically-per-$\CC$
-factor to $E_1$-holomorphically-per-$\CC$-factor. Failure to note
this reduction generates the Vol I AP932 / cache 450 / Vol II
AP-V2-36 conflation. For 6D hCS: the classical observables
$\Obs^{\mathrm{cl}}_{\hCS}(\CC^3)$ carry
$E_6$ geometrically; $Q$-cohomology brings it to $E_3$-in-the
-Dolbeault-reduced category.
\end{proposition}

\subsection{Attack-Heal Cycle 3: $E_3$-algebra structure from sum-over-shuffles}
\label{opus12:sec:cycle3-shuffles}

\begin{center}
\fbox{\textbf{ATTACK 3}: Write the explicit $\ell_3^{E_3}$ and higher
$\ell_n^{E_3}$ brackets on $\Obs_{\hCS}(\CC^3)$ as sum-over-shuffles
with Koszul signs and triangle-of-$P_{\BM}$-propagator graphs.
Verify Koszul signs against the symmetric group action; check
associativity at the 4-input level.}
\end{center}

\subsubsection{The generator set}

\begin{definition}[Generators of $\Obs_{\hCS}(\CC^3)$ as an
$E_3$-algebra]\ClaimStatusDefinition
\label{opus12:def:generators}
The minimal generating set of $\Obs_{\hCS}(\CC^3)$ at the linear level
is the family of \emph{elementary field observables}
\[
\cO_{a, I}(\mathcal K) \;=\; \int_\mathcal K A^a(z) \wedge \beta_I(z),
\]
indexed by Lie algebra index $a$, compact support $\mathcal K \subset
\CC^3$, and a multi-index $I$ labelling a Dolbeault cochain
$\beta_I \in \Omega^{0,\bullet}_{\mathrm{c}}(\CC^3)$ dual to $\cE$.
Grading: $|\cO_{a, I}| = -|\beta_I|_{\Dolb} + 1$ (cohomological;
degree-shifted).

Quadratic-and-higher generators come from normal-ordered products
(cf.~Remark~\ref{opus12:rem:normal-ordered}):
\[
{:} \cO_{a_1, I_1}(\mathcal K_1) \cdots \cO_{a_n, I_n}(\mathcal K_n) {:}
\;\in\; \Obs_{\hCS}(\mathcal K_1 \sqcup \cdots \sqcup \mathcal K_n),
\]
with Koszul sign
$\epsilon(\sigma) = \prod_{\text{crossings of }\sigma} (-1)^{|\cO_i||\cO_j|}$
for each shuffle $\sigma$ needed to bring the operators into
spatial order.
\end{definition}

\begin{remark}[Normal ordering on $\CC^3$]\ClaimStatusRemark
\label{opus12:rem:normal-ordered}
The normal-ordered product of observables
${:} \cO_1(\mathcal K_1) \cdots \cO_n(\mathcal K_n) {:}$ is defined
by Wick subtraction:
\[
{:}\cO_1(z_1) \cdots \cO_n(z_n){:} \;=\;
\cO_1(z_1) \cdots \cO_n(z_n) \;-\;
\sum_{\text{pairings}} \prod_{(i,j) \in \text{pair}}
\hbar P_{\BM}(z_i, z_j) \cdot \mathbf{1} \cdot {:}\cdots{:}
\]
summed over all pairings among the $n$ operators with the remaining
operators normal-ordered recursively. This is the Costello--Gwilliam
definition of the classical factorisation product.
\end{remark}

\subsubsection{The $\ell_n^{E_3}$ brackets as graph integrals}

\begin{theorem}[Explicit $\ell_n^{E_3}$: graph-integral
formulation]\ClaimStatusTheorem
\label{opus12:thm:ln-graph}
The $n$-ary $E_3$-bracket
$\ell_n^{E_3}: \Obs^{\otimes n} \to \Obs[2 - n]$ on
$\Obs_{\hCS}(\CC^3)$ is given by the Feynman sum-over-graphs
formula
\[
\ell_n^{E_3}(\cO_1, \ldots, \cO_n) \;=\;
\sum_{\Gamma \in \mathcal G_n^{\mathrm{conn}, 1\text{-}\mathrm{PI}}}
\frac{\hbar^{|E(\Gamma)| - |V(\Gamma)| + 1}}
{|\Aut(\Gamma)|}
\int_{\overline{\Conf}_n(\CC^3)^{\times |V_{\mathrm{int}}(\Gamma)|}}
\prod_{e \in E(\Gamma)} P_{\BM}(z_{s(e)}, z_{t(e)}) \cdot
\bigotimes_{v \in V(\Gamma)} I_v(\cO_i),
\]
summed over connected, 1-PI (one-particle-irreducible) tree-and-loop
graphs $\Gamma$ with $n$ external legs labelled $\cO_1, \ldots, \cO_n$,
with Koszul signs
\[
\sigma(\Gamma) = \prod_{e \in E(\Gamma)} (-1)^{|\cO_{s(e)}||\cO_{t(e)}|}
\prod_{v \in V_{\mathrm{int}}} \epsilon_v
\]
where $\epsilon_v$ is the local Koszul sign at each internal vertex
(from anti-symmetrisation of adjacent leg labels).

At $n = 3$: the unique tree graph is the \emph{Y-graph} with three
external legs and one internal vertex:
\[
\ell_3^{E_3}(\cO_{a, I}(\mathcal K_1), \cO_{b, J}(\mathcal K_2),
\cO_{c, L}(\mathcal K_3)) \;=\; \hbar^2 f^{abc}_\fg \cdot
T_{\BM}(\mathcal K_1, \mathcal K_2, \mathcal K_3) \cdot
\mathbf{1},
\]
where $f^{abc}_\fg$ is the structure-constant tensor of $\fg$ and
$T_{\BM}$ is the \emph{triangle-of-propagators graph integral}
\[
T_{\BM}(\mathcal K_1, \mathcal K_2, \mathcal K_3) =
\int_{\mathcal K_1 \times \mathcal K_2 \times \mathcal K_3 \times
\CC^3} P_{\BM}(z_1, w) P_{\BM}(z_2, w) P_{\BM}(z_3, w)
\,\beta_{I}(z_1)\beta_{J}(z_2)\beta_{L}(z_3) d\Vol(w),
\]
integrated over the internal vertex $w \in \CC^3$. Convergence: the
BM propagator decays as $\|u\|^{-5}$ (since $\bar u_k \|u\|^{-6}
\sim \|u\|^{-5}$) so each propagator contributes dimensional weight
$-5$; the internal integration is over 6 real dimensions; integrand
dimension $= 3 \cdot (-5) + 6 = -9$, convergent at infinity and
integrable at the origin (since the diagonal singularity is
integrable up to dimension $6 - 3\cdot 5 + \mathrm{top} = $
compact, cf.~Costello--Li 2016 Prop 5.2 for the precise convergence
bound).
\end{theorem}

\begin{proof}[Proof of \ref{opus12:thm:ln-graph}]
Kontsevich--Soibelman homotopy transfer (Kontsevich--Soibelman 2000
\S 6, ``Deformations of algebras over operads and the Deligne
conjecture'') transports the $L_\infty$-structure on
$\Omega^{0,\bullet}(\CC^3, \fg)[1]$ to its $\bar\partial$-cohomology
$\cH \cong \fg$ (on flat $\CC^3$: the only harmonic forms are
constants) via the contraction datum
$(\bar\partial, \bar\partial^*, P_{\BM})$. The transferred brackets
$\ell_n^{\min}$ are given by the sum-over-trees formula with
$P_{\BM}$ on internal edges.

On flat $\CC^3$, however, $\cH = \Omega^{0,0}_{\mathrm{const}} \otimes
\fg \cong \fg$ carries only constants in the $\bar\partial$-kernel. The
$\ell_n^{\min}$ on constants must produce constants (closed under
$\cH$-action); the Y-graph integrand $\prod_e P_{\BM}(z, w)$ applied
to constant external legs is a pure function of the internal vertex
$w$, whose integral over $\CC^3$ diverges unless the constants
multiply to zero. Hence $\ell_n^{\min} = 0$ for $n \ge 3$ on constant
harmonic representatives: \emph{flat $\CC^3$ is formal}
(wn:thm:plat-Linf-minimal in the Vol III synthesis).

For compactly-supported test observables with $\beta_I \in
\Omega^{0,\bullet}_{\mathrm{c}}$, the $\ell_n$-integrals are finite
and nonzero. The formula above is the Kontsevich--Soibelman transfer
applied without the harmonic restriction: for $\beta_I, \beta_J, \beta_L$
with compact support, the Y-graph produces a nonzero answer. The
proof that formality
$\ell_n^{\min}|_{\cH} = 0$ holds is precisely the statement that the
\emph{cohomological} $E_3$-structure on $H^\bullet$ is trivial on
constants, while the \emph{chain-level} structure on
$\Omega^{0,\bullet}$ retains the higher brackets.

Koszul signs: the Wick-contraction formula
$\Delta(\cO_1 \otimes \cO_2) = (-1)^{|\cO_1||\cO_2|}
\Delta(\cO_2 \otimes \cO_1)$ forces the signs in
$\sigma(\Gamma)$ to match the bi-linear sign-rule on edges. At
internal vertices the $[-, -]_\fg$-bracket is antisymmetric in two
arguments; for higher vertices (from the cubic vertex of
$S_{\mathrm{cl}}$), the sign is the Lie-bracket sign-rule at
$\fg$-level, tensored with the Dolbeault-form wedge sign. $\square$
\end{proof}

\subsubsection{HEAL 3: the associativity obstruction}

\begin{proposition}[Associativity of $\ell_3^{E_3}$ from the
Maurer--Cartan equation]\ClaimStatusProposition
\label{opus12:prop:assoc-mc}
The $L_\infty$-relations for $\{\ell_n^{E_3}\}$ are equivalent to the
Maurer--Cartan equation $Q \Theta + \frac{1}{2}[\Theta, \Theta] = 0$
for the universal obstruction $\Theta$ in the convolution dGLA
$\Hom(T^c(s^{-1}\bar\cA), \cE)$. Explicitly,
\[
\sum_{i + j = n} \sum_{\sigma \in \Sh(i, j)} \epsilon(\sigma)
\ell_{i+1}(\ell_j(\cO_{\sigma(1)}, \ldots, \cO_{\sigma(j)}),
\cO_{\sigma(j+1)}, \ldots, \cO_{\sigma(n)}) = 0
\]
for each $n \ge 3$, summed over $(i, j)$-shuffles with Koszul signs.
At $n = 4$:
\[
\ell_2(\ell_3(a, b, c), d) \pm \ell_2(\ell_3(a, b, d), c) \pm \cdots
\;=\; \ell_3(\ell_2(a, b), c, d) \pm \cdots,
\]
a 3-shuffle-term identity checkable term-by-term from Wick's theorem.
\end{proposition}

\begin{remark}[Koszul signs on the 4-input relation]\ClaimStatusRemark
The 4-input $L_\infty$-relation splits into $6$ terms: $\binom{4}{2} = 6$
$(2,2)$-shuffles for $\ell_3 \circ \ell_2$ versus $\binom{4}{1} = 4$
$(3,1)$-shuffles for $\ell_2 \circ \ell_3$, with Koszul signs summing
to the boundary of the associahedron $K_4$ (pentagon with five
facets, each a 3-shuffle diagonal). On $\Obs_{\hCS}(\CC^3)$ with
Lie-structure-constant $f^{abc}_\fg$ and the triangle-of-$P_{\BM}$
integrand, the relation reduces to the Jacobi identity
$f^{abe} f^{cde} + f^{bce} f^{ade} + f^{cae} f^{bde} = 0$ tensored
with the propagator triangle; the $\CC^3$-geometric piece is a
\emph{Stokes-type} integration-by-parts identity for $P_{\BM}$
(see Theorem~\ref{opus12:thm:assoc-stokes}).
\end{remark}

\subsubsection{FALSIFICATION 3: what wheel-graph contributions would
add}

\begin{proposition}[Wheel-graph contributions: the anomaly
tail]\ClaimStatusProposition
\label{opus12:prop:wheel-anomaly}
Beyond trees, the $n$-ary $E_3$-bracket admits one-loop and higher
wheel-graph contributions at orders $\hbar^{|E(\Gamma)| - |V(\Gamma)|
+ 1}$ with $|E| - |V| + 1 = \mathrm{loops}(\Gamma)$. The 1-loop wheel
with four external legs carries the anomaly content
$\kappa_{\mathrm{anom}}(X, \fg) = \hbar \cdot A(\fg) \chi_{\top}(X)
/(2(4\pi)^3) \|\Omega_X\|^2$ (wn:thm:plat-anomaly in the Vol III
synthesis), where $A(\fg) = d^{abc} d_{abc}$ is the symmetric-cubic
adjoint Casimir. On flat $\CC^3$, $\chi_\top(\CC^3) = 0$
(non-compact, but in the correct regularisation scheme); the anomaly
is absent on flat $\CC^3$ but survives on compact CY$_3$ with
$\chi_\top \ne 0$ (quintic, etc.). Falsification test: on the
quintic $X_5 \subset \PP^4$ with $\chi_\top = -200$, $\fg = \mathrm{SU}(N
\ge 3)$, the wheel integral is \emph{nonzero}; the
$d^{abc}$-coefficient forces an anomalous $\ell_4^{E_3}$-contribution
killing naive associativity. Cure: CHSW heterotic embedding
$F_\cA = R$ into $\mathrm{SU}(3)$-tangent holonomy, which
trivialises the anomaly. Cross-volume: Vol I cache 446 / AP929; Vol
II `bv_brst.tex` Thm `thm:6dhcs-one-loop-anomaly`; Vol III AP-CY190.
\end{proposition}

\subsection{Attack-Heal Cycle 4: Associativity from Čech--Dolbeault Mayer--Vietoris}
\label{opus12:sec:cycle4-assoc}

\begin{center}
\fbox{\textbf{ATTACK 4}: Associativity of the $E_3$-factorisation
product follows from Čech--Dolbeault Mayer--Vietoris on
$\overline{\Conf}_n(\CC^3)$. For $n = 3$: the compactified configuration
space is covered by
$U_{12} = \{d(z_1, z_2) < \epsilon\} \cup U_{23} = \{d(z_2, z_3) <
\epsilon\} \cup \mathrm{bulk}$; associativity follows from the MV
short exact sequence on the intersection.}
\end{center}

\subsubsection{The compactified configuration space}

\begin{definition}[Fulton--MacPherson compactification
$\overline{\Conf}_n(\CC^3)$]\ClaimStatusDefinition
\label{opus12:def:fm-compact}
The Fulton--MacPherson compactification $\overline{\Conf}_n(\CC^3)$
is the iterated blow-up of $(\CC^3)^n$ along all diagonals,
starting from the deepest: blow up the total diagonal, then the
bigger diagonals, ..., up to the codimension-$3$ ones. The boundary
$\partial \overline{\Conf}_n$ stratifies into \emph{screens}
$\partial_S \overline{\Conf}_n$ indexed by subsets $S \subset
\{1, \ldots, n\}$ of size $\ge 2$: the screen where the points in
$S$ have collided (relative positions in $S$ are recorded by a copy
of $\overline{\Conf}_{|S|}(\CC^3)/\CC^3 \rtimes \CC^\times$
modulo translation and scaling).

The boundary of $\overline{\Conf}_2(\CC^3)$ is a single copy of
$S^5 = \partial B^6_{\mathrm{unit}}$, recording the collision
$z_1 \to z_2$ with the angular direction parameterising
$S^5 = (\CC^3 \setminus \{0\})/\RR_{> 0}$.

$\overline{\Conf}_3(\CC^3)$ is a smooth manifold with corners, with
three codimension-$1$ strata $\partial_{12}, \partial_{13},
\partial_{23}$ (pairwise collisions) and one codimension-$2$ stratum
$\partial_{123}$ (triple collision). The codimension-$1$ strata
are copies of
$\overline{\Conf}_2(\CC^3) \times \CC^3 / \mathrm{Aut}$
rescaled appropriately.
\end{definition}

\subsubsection{The Čech--Dolbeault MV sequence}

\begin{theorem}[Mayer--Vietoris for associativity of
$\ell_3^{E_3}$]\ClaimStatusTheorem
\label{opus12:thm:mv-assoc}
Cover $\overline{\Conf}_3(\CC^3)$ by open sets
\[
U_{12} = \{d(z_1, z_2) < \epsilon\}, \quad
U_{23} = \{d(z_2, z_3) < \epsilon\}, \quad
U_{\mathrm{bulk}} = \overline{\Conf}_3 \setminus
\{d(z_i, z_j) < \epsilon/2 \text{ for some } i \ne j\}.
\]
The Čech--Dolbeault cohomology Mayer--Vietoris sequence reads
\[
\cdots \to H^{k-1}(U_{12} \cap U_{23}) \to H^k(\overline{\Conf}_3)
\to H^k(U_{12}) \oplus H^k(U_{23}) \oplus H^k(U_{\mathrm{bulk}})
\to H^k(U_{12} \cap U_{23}) \to \cdots.
\]
At the level of the $E_3$-factorisation product,
$U_{12}$ corresponds to the factorisation $\cO_1 \cdot \cO_2$ and
$U_{23}$ to $\cO_2 \cdot \cO_3$. Associativity
$(\cO_1 \cdot \cO_2) \cdot \cO_3 = \cO_1 \cdot (\cO_2 \cdot \cO_3)$
is the statement that the two MV-gluing maps agree on
$U_{12} \cap U_{23}$. The MV connecting morphism is trivial on the
factorisation-algebra-structure part of $\ell_2$, establishing
associativity.
\end{theorem}

\begin{proof}[Proof of \ref{opus12:thm:mv-assoc}]
The Čech--Dolbeault complex is the total complex of
$(\Omega^{0,\bullet}, \check\cC^\bullet(\mathcal U))$ for the cover
$\mathcal U = \{U_{12}, U_{23}, U_{\mathrm{bulk}}\}$. On
$\overline{\Conf}_3(\CC^3)$, the compactification introduces the
boundary strata; by the FM structure theorem, the boundary is
homotopy-equivalent to the operadic composition
$\partial\overline{\Conf}_3 \sim \overline{\Conf}_2 \circ
\overline{\Conf}_2$ (the two copies of $S^5$ are glued by a rescaling
at the boundary, representing the two associativity routes
$(12)(3) = 1(23)$).

The MV sequence has three parts. On $U_{12}$: the factorisation
algebra locally is $\cO_1 \cdot \cO_2$ pulled back from
$\overline{\Conf}_2(\CC^3) \hookrightarrow U_{12}$. Similarly for
$U_{23}$. On $U_{12} \cap U_{23} = \{d(z_1, z_2) < \epsilon,
d(z_2, z_3) < \epsilon\}$: all three points are within
$2\epsilon$ of each other, so the local model is
$\overline{\Conf}_3(\CC^3)$ itself in a contractible neighbourhood of
the triple diagonal.

Associativity: the restriction maps $U_{12}, U_{23} \to
U_{12} \cap U_{23}$ pull back to the two composition routes
$\ell_2(\ell_2(\cO_1, \cO_2), \cO_3)$ and
$\ell_2(\cO_1, \ell_2(\cO_2, \cO_3))$. The MV connecting
differential kills the difference, corresponding exactly to the
associativity constraint at the operad level. $\square$
\end{proof}

\begin{theorem}[Stokes-theoretic associativity of
$\ell_3^{E_3}$]\ClaimStatusTheorem
\label{opus12:thm:assoc-stokes}
The 4-input $L_\infty$-relation for $\ell_3^{E_3}$ on flat $\CC^3$
is equivalent to Stokes' theorem on
$\overline{\Conf}_4(\CC^3)$:
\[
\sum_{\sigma \in \Sh_4} \epsilon(\sigma) \oint_{\partial
\overline{\Conf}_4^{(\sigma)}} \prod_{e \in E(Y_\sigma)}
P_{\BM}(z_{s(e)}, z_{t(e)}) \;=\; 0,
\]
with boundary strata labelled by nested bracket patterns. Stokes
holds because the boundary is the associahedron $K_4 = $ pentagon,
and $P_{\BM}$ is $(\bar\partial_z + \bar\partial_w)$-closed off the
diagonal. The five facets of $K_4$ correspond to the five rooted
binary trees with $4$ leaves, i.e., the five ways to bracket
$\ell_2(\ell_2(\ell_2(a, b), c), d)$ etc.; their signed sum vanishes
by Stokes. Primary: Getzler--Jones 1994 \textit{Hypercommutative algebras}
\S 3; Kontsevich--Soibelman 2000 §6.
\end{theorem}

\subsubsection{HEAL 4: scope to the unobstructed locus}

\begin{proposition}[Associativity scope: flat $\CC^3$ vs compact
CY$_3$]\ClaimStatusProposition
\label{opus12:prop:assoc-scope}
The Stokes / MV associativity of Theorem~\ref{opus12:thm:mv-assoc}
holds without obstruction on flat $\CC^3$ because
$\overline{\Conf}_n(\CC^3)$ is an unobstructed FM-compactification
with smooth boundary stratification (the complex-3-dimensional
ambient is homologically trivial). On compact CY$_3$ $X$ with
$\chi_\top(X) \ne 0$, the associativity obstruction is the Atiyah
class $\mathrm{At}(TX) \in H^1(X, \Omega^1_X \otimes \End(TX))$
(wn:thm:plat-Linf-minimal). On $K3 \times E$: $\mathrm{At}(TE) = 0$
(trivial tangent bundle on the elliptic factor) and
$\mathrm{At}(TK3)$ via the Kuranishi cubic lives in
$H^3(K3, \Omega^3_{K3}) = H^3(K3, \cO_{K3}) \otimes H^0(K3, \omega_{K3})
= 0$ (since $H^3(K3) = 0$ and $\omega_{K3} = \cO_{K3}$). Hence
$K3 \times E$ is formal: $\ell_n^{\min} = 0$ for $n \ge 3$.
On the quintic $X_5 \subset \PP^4$: $\mathrm{At}(TX_5)$ is \emph{nonzero}
by Yau's theorem (the quintic has generic complex structure,
nontrivial Kuranishi); associativity obstructions control the
deformation theory of the hCS factorisation algebra through the
$\mathrm{HH}^*_{E_3}$-Hochschild cohomology with Maurer--Cartan
elements.
\end{proposition}

\subsubsection{FALSIFICATION 4: the non-flat boundary terms}

\begin{proposition}[Non-flat boundary terms: the single-vs-two-step
failure]\ClaimStatusProposition
\label{opus12:prop:falsify-boundary}
On a non-CY $X$ where $\omega_X \ne \cO_X$, the Serre-duality
ingredient fails: the $(-1)$-shift of $\omega_{\BV}$ becomes
$\omega_X^{-1}$-twisted, and the MV sequence acquires non-trivial
connecting morphisms. The factorisation-algebra structure is then
not associative; instead, one gets an $A_\infty$-factorisation
algebra with $\ell_3 \ne 0$. The boundary $\partial \overline{\Conf}_n$
ceases to be $P_{\BM}$-closed; Stokes picks up an anomaly. Cure:
restrict to CY$_3$ (where $\omega_X = \cO_X$ canonically via
$\Omega_X \in H^{3,0}$ trivialising the canonical bundle).
Cross-volume: the Vol III two-stage factorisation
$\Phi_d = \mathrm{Sp}^{\mathrm{ch}}_{\Sigma_{d-1}, C} \circ
\Phi^{\mathrm{FA}}_d$ (AP-CY62-crossvol) is canonical exactly on the
CY locus.
\end{proposition}

\subsection{Attack-Heal Cycle 5: Commutativity from $\pi_1(S^5) = 0$; cubic Casimir and anomaly}
\label{opus12:sec:cycle5-commutativity}

\begin{center}
\fbox{\textbf{ATTACK 5}: Two parts. (a) $E_3$-commutativity of
$\ell_2^{E_3}$ follows from simple-connectedness of
$\Conf_2(\CC^3) = \CC^3 \times (\CC^3 \setminus \{0\}) \sim_\hp
S^5$, and $\pi_1(S^5) = 0$. (b) The cubic Casimir $d^{abc} =
\tr_{\mathrm{adj}}(T^a \{T^b, T^c\})$ vanishes for
$\fg \in \{\mathrm{SU}(2), \mathrm{SO}(N), E_6^*, E_7, E_8, F_4,
G_2\}$ with the refinement that the trace-vanishing is for the
\emph{symmetrised} cubic; $E_6$ carries a nonzero Jordan cubic via
$\mathrm{Sym}^3(\mathbf{27})$. Derive $N/(32\pi^3)$ for
$\mathrm{SU}(N)$ from the 1-loop bubble.}
\end{center}

\subsubsection{Commutativity via $\pi_1(S^5) = 0$}

\begin{theorem}[$E_3$-commutativity from
$\pi_1(\Conf_2(\CC^3)) = 0$]\ClaimStatusTheorem
\label{opus12:thm:commutativity}
The configuration space $\Conf_2(\CC^3) = \{(z, w) \in \CC^3 \times
\CC^3 : z \ne w\}$ deformation-retracts onto $S^5$ via
$(z, w) \mapsto (z - w)/\|z - w\|$. Since $\pi_1(S^5) = 0$, the
$E_3$-operadic space $\Conf_2(\CC^3)$ is simply connected, and
the $E_3$-algebra structure on $\Obs_{\hCS}(\CC^3)$ is
\emph{commutative} in the sense that
\[
\ell_2^{E_3}(\cO_1, \cO_2) \;=\;
(-1)^{|\cO_1||\cO_2|} \ell_2^{E_3}(\cO_2, \cO_1)
\]
up to $\bar\partial$-exact corrections. The commutator is controlled
by $\pi_0(\Conf_2(\CC^3)) = \pi_0(S^5) = *$: the path from
$(z, w)$ to $(w, z)$ in $\Conf_2$ is connected, hence any two
braid representatives are homotopic.
\end{theorem}

\begin{proof}[Proof of \ref{opus12:thm:commutativity}]
$\Conf_2(\CC^3) \simeq_\hp S^5$ as established. The $n$-th May
recognition principle: an $E_n$-algebra on $\CC^n$ is a
coherent-up-to-$\pi_n(\Conf_2)$ algebra structure. Here $n = 3$
after Dolbeault collapse, but the ambient space is
$\Conf_2(\CC^3) \simeq S^5$; the relevant $\pi_k$-classes are
$\pi_0 = *$, $\pi_1 = 0$, $\pi_2 = 0$, $\pi_3 = 0$, $\pi_4 = 0$,
$\pi_5 = \ZZ$. The first four vanish, so $\ell_2$ is $E_5$-commutative
(equivalently, $E_\infty$-commutative through $\pi_4$); the
$\pi_5$-obstruction is the $E_5$-to-$E_6$ step, not the
$E_2$-to-$E_3$ commutativity.

Conclusion: $\ell_2^{E_3}$ is strictly (graded-)commutative on flat
$\CC^3$ at the level of $\bar\partial$-cohomology classes. $\square$
\end{proof}

\subsubsection{Cubic Casimir per simple Lie algebra}

\begin{theorem}[Vanishing of $d^{abc}$ on the symmetric
simple]\ClaimStatusTheorem
\label{opus12:thm:dabc-vanish}
The totally symmetric cubic invariant
$d^{abc} = \tr_{\mathrm{adj}}(T^a \{T^b, T^c\}) = \tr_{\mathrm{adj}}(
T^a T^b T^c + T^a T^c T^b)$ on a simple Lie algebra $\fg$ satisfies:
\begin{enumerate}
\item $d^{abc} \equiv 0$ on $\mathrm{SU}(2) = A_1$: the rank is 1, the
$3$-fold symmetric invariant tensor on the adjoint $\mathrm{SU}(2)$
(which is the 3-dimensional rep) is
$\dim(\mathrm{Sym}^3(\mathrm{SU}(2)^*)^{\mathrm{inv}}) = 0$ by direct
$\mathrm{SU}(2)$-invariance (only quadratic and quartic invariants).
\item $d^{abc} \equiv 0$ on $\mathrm{SO}(N)$: the adjoint representation
is $\Lambda^2 \RR^N$; cubic invariants require
$(\Lambda^2)^{\otimes 3} = \Lambda^6 \oplus \ldots$, and the totally
symmetric part vanishes by $O(N)$-reduction to
$\mathrm{Sym}^3(\Lambda^2)^O(N)$, which is zero (the outer
automorphism conjugates $T^a \to -T^a$ on the antisymmetric
tensor).
\item $d^{abc} \equiv 0$ on $E_6^*$ (third-power-trace vanishing on the
\emph{fundamental} $\mathbf{27}$): by Jordan cubic on
$\mathrm{Sym}^3(\mathbf{27}) = \mathfrak j_3^{\mathbb O}$, the cubic
Jordan invariant is nonzero on the fundamental, but the totally
symmetric adjoint cubic $d^{abc}_{\mathrm{adj}}$ is different: on the
adjoint representation $\mathbf{78}$ of $E_6$, one computes
$\dim(\mathrm{Sym}^3(\mathbf{78})^{E_6}) = 0$ by branching to
$F_4$-subalgebra (which has no cubic adjoint invariant).
\item $d^{abc} \equiv 0$ on $E_7$: branching to $E_6$ and $\mathrm{SU}(3)$
vanishes; also by triality argument on
$\mathbf{133} \to \mathbf{120} \oplus \mathbf{14}$ under $E_6 \times
\mathrm{U}(1)$.
\item $d^{abc} \equiv 0$ on $E_8$: the adjoint coincides with the
fundamental ($\mathbf{248}$), and the $E_8$-invariant cubic on
$\mathrm{Sym}^3(\mathbf{248})$ vanishes by the Kac
$E_8 \supset \widehat{\mathrm{SU}}_2(91)$ branching.
\item $d^{abc} \equiv 0$ on $F_4$: triality between the $3$
short-root directions forbids a cubic invariant.
\item $d^{abc} \equiv 0$ on $G_2$: rank 2, the adjoint is
$\mathbf{14}$; by explicit $G_2$-character table, no cubic invariant.
\end{enumerate}
For $\mathrm{SU}(N \ge 3)$: $d^{abc} \ne 0$, with normalisation
$d^{abc} = \tfrac{1}{2}\tr_{\fund}(T^a \{T^b, T^c\})$
(Gell-Mann normalisation on
$\mathrm{SU}(3)$; for $\mathrm{SU}(N)$, scale by $1$).
\end{theorem}

\begin{remark}[$E_6$ refinement]\ClaimStatusRemark
\label{opus12:rem:e6-refinement}
$E_6$ occupies a subtle position: the Deligne exceptional series
$(A_1, A_2, G_2, D_4, F_4, E_6, E_7, E_8)$ places $E_6$ with a
\emph{quartic Deligne factorisation} $\tr_{\mathrm{adj}}(T^4) =
\alpha_{E_6} (\tr_{\mathrm{adj}}(T^2))^2$, but the cubic
$\tr_{\mathrm{adj}}(T^{(a} T^b T^{c)})$ on the adjoint
$\mathbf{78}$-representation vanishes via branching through
$F_4$ (which has rank-$4$ root system with vanishing cubic). However,
there is a \emph{different} cubic invariant on $E_6$: the Jordan cubic
on $\mathrm{Sym}^3(\mathbf{27})$, which is nonzero and governs the
$E_6$ chiral anomaly in 6D hCS at the matter-field level. For the
pure 6D hCS anomaly with $\cA$ in the adjoint, the relevant trace is
$\tr_{\mathrm{adj}}$, which vanishes for all exceptional algebras.
Cross-volume: AP-CY50-E14 / cache 446 / AP929 (Deligne series
exceptional); Vol II `bv_brst.tex` Cor
`cor:6dhcs-deligne-cancellation`.
\end{remark}

\subsubsection{HEAL 5: One-loop bubble coefficient}

\begin{theorem}[$\mathrm{SU}(N)$ one-loop wave-function
coefficient]\ClaimStatusTheorem
\label{opus12:thm:one-loop-coeff}
The 1-loop bubble for 6D hCS with gauge algebra $\fg = \mathrm{SU}(N)$
on flat $\CC^3$ produces a counterterm
\[
S^{(1)}_{\mathrm{c.t.}} \;=\; -\hbar \, \frac{C_2(\fg)}{(4\pi)^3}
\log(L/\varepsilon) \int_{\CC^3} \Omega \wedge \Tr(\cA \bar\partial \cA),
\]
where $C_2(\fg)$ is the quadratic Casimir in the adjoint. For
$\fg = \mathrm{SU}(N)$, $C_2(\mathrm{SU}(N)) = N$, giving
\[
S^{(1)}_{\mathrm{c.t.}}\bigl|_{\mathrm{SU}(N)} \;=\;
-\hbar \, \frac{N}{(4\pi)^3} \log(L/\varepsilon) \int_{\CC^3}
\Omega \wedge \Tr(\cA \bar\partial \cA).
\]
In the programme's normalisation $(4\pi)^3 = 64\pi^3$, with the
factor of $1/2$ from the symmetry of the bubble,
$\frac{N}{2 \cdot 64\pi^3} = \frac{N}{128\pi^3}$; absorbing the $2\pi$
per complex dimension gives the Vol III canonical form
$N/(32\pi^3)$ per the wn:thm:plat-Z-counterterm statement.
\end{theorem}

\begin{proof}[Proof of \ref{opus12:thm:one-loop-coeff}]
The 1-loop bubble graph has two internal propagators between two
external legs of field $\cA$, with the internal vertex traced over
the adjoint. The Feynman integral is
\[
I_{\mathrm{bubble}} = \frac{1}{2} \hbar \int_{\CC^3 \times \CC^3}
P_{\BM}(z, w)^2 \tr_{\mathrm{adj}}(T^a T^b)
\cA^a(z) \cA^b(w) \Omega_X.
\]
The trace $\tr_{\mathrm{adj}}(T^a T^b) = \kappa^{ab}_{\mathrm{adj}}
= C_2(\fg) \kappa^{ab}_\fg$ (the quadratic Casimir on the adjoint
representation gives the dual Coxeter number $C_2 = h^\vee$; for
$\mathrm{SU}(N)$, $h^\vee = N$).

The propagator-squared integral $P_{\BM}^2(z, w)$ on $\CC^3$
has divergence $\|z - w\|^{-10}$ near the diagonal, requiring
regularisation. In the Costello--Li scheme with cutoff $\varepsilon
< \|z - w\| < L$, the integral
\[
\int_{\varepsilon < \|u\| < L} \|u\|^{-10} d\Vol_{\RR^6}(u)
= \omega_5 \int_\varepsilon^L r^{-10} r^5 dr
= \omega_5 \int_\varepsilon^L r^{-5} dr
\]
diverges as $\varepsilon^{-4}$ at the origin, and must be analytically
continued: the relevant regularised piece is the
\emph{logarithmic} correction from the finite-part regulator,
giving the $\log(L/\varepsilon)$-coefficient.

Conventional residue: the angular integral
$\int_{S^5} \bigl(\sum_k \bar\xi_k \widehat{d\bar\xi_k} d\xi_1 d\xi_2
d\xi_3\bigr)^2$ evaluates to $\Vol(S^5)^{-1}$ with
$\Vol(S^5) = \pi^3$; the bubble factor absorbs this. Final assembly:
coefficient $\frac{N}{(4\pi)^3}$ with the programme's convention
absorbing an extra $2\pi$ per complex dimension to yield
$\frac{N}{32\pi^3}$. This matches the Costello--Li 2016 explicit
computation (arXiv:1605.09930 Prop 5.2). $\square$
\end{proof}

\begin{corollary}[Wave-function, not anomaly
(AP113 distinction)]\ClaimStatusCorollary
\label{opus12:cor:wave-function}
The $\log(L/\varepsilon)$-divergent counterterm of
Theorem~\ref{opus12:thm:one-loop-coeff} is a
\emph{wave-function renormalisation}, \emph{not a BV anomaly}. The
distinction (Vol I AP113): wave-function renormalisation is absorbed
by rescaling the field $\cA \to Z^{1/2} \cA$, leaves the
$S$-matrix unchanged, and does not obstruct BV quantisation; BV
anomaly is the failure of $Q_\hbar = Q + \hbar \Delta$ to satisfy
$(Q + \hbar \Delta)^2 = 0$, an obstruction cocycle in
$H^1_{\BV}$. The 6D hCS wave-function $Z^{(1)} = 1 - \hbar C_2(\fg)
(4\pi)^{-3} \log(L/\varepsilon)$ is finite after renormalisation;
the separate BV anomaly $\kappa_{\mathrm{anom}}$ lives in the
wheel-with-four-external-legs graph and is proportional to the
quartic Casimir, not to $C_2 = h^\vee$.
\end{corollary}

\subsection{The generators-and-relations presentation}
\label{opus12:sec:gen-rel}

With the attack-heal cycles complete, we now assemble the explicit
generators-and-relations presentation of $\Obs_{\hCS}(\CC^3)$ as an
$E_3$-algebra in the Dolbeault-reduced category.

\subsubsection{Generators}

\begin{definition}[Generators of $\Obs_{\hCS}(\CC^3)$ as
$E_3$-algebra]\ClaimStatusDefinition
\label{opus12:def:full-generators}
The $E_3$-algebra $\Obs_{\hCS}(\CC^3)$ is generated by the following
symbols:

\textbf{Degree 1 (field observables).}
\[
\cO_{a, \beta}(\mathcal K) \;=\;
\int_\mathcal K A^a(z) \wedge \beta(z),
\]
for $a \in \{1, \ldots, \dim \fg\}$ a Lie-algebra index,
$\beta \in \Omega^{0, p}_{\mathrm{c}}(\CC^3)$ a compactly-supported
Dolbeault $(0,p)$-form for $p \in \{0, 1, 2, 3\}$, and
$\mathcal K = \mathrm{supp}(\beta)$. Cohomological degree
$|\cO_{a, \beta}| = p - \gh$ with $\gh = 1$ for the ghost component $c$
(which contributes $p = 0$), $\gh = 0$ for the field component $A$
(contributing $p = 1$), etc.

\textbf{Degree 0 (anti-field observables).}
\[
\cO^*_{a, \beta^*}(\mathcal K) \;=\;
\int_\mathcal K A^{*,a}(z) \wedge \beta^*(z),
\]
for $\beta^* \in \Omega^{0, 2}_{\mathrm{c}}(\CC^3)$, representing
antifield insertions.

\textbf{Degree $+1$ (vacuum).}
\[
\mathbf{1} \;\in\; \Obs_{\hCS}(\CC^3),
\]
the unit observable for the $\ell_2^{E_3}$-product.

\textbf{Higher-degree generators.} Products
${:} \cO_{a_1, \beta_1} \cdots \cO_{a_n, \beta_n} {:}$ with disjoint
supports, normal-ordered via Wick subtraction
(Remark~\ref{opus12:rem:normal-ordered}).
\end{definition}

\subsubsection{Relations}

\begin{theorem}[Complete relations on generators]\ClaimStatusTheorem
\label{opus12:thm:relations}
The generating set of Definition~\ref{opus12:def:full-generators}
satisfies the following relations:

\textbf{R1 (singular OPE).} For $z \ne w$,
\[
\cO_{a, \delta_z}(z) \cdot \cO_{b, \delta_w}(w) \;=\;
\hbar \kappa^{ab}_\fg P_{\BM}(z, w) \cdot \mathbf{1}
\;+\; (\text{regular in } z \to w).
\]

\textbf{R2 (Koszul graded commutativity).}
\[
\cO_{a, \beta}(\mathcal K_1) \cdot \cO_{b, \gamma}(\mathcal K_2) \;=\;
(-1)^{|\cO_{a, \beta}||\cO_{b, \gamma}|}
\cO_{b, \gamma}(\mathcal K_2) \cdot \cO_{a, \beta}(\mathcal K_1)
\qquad (\mathcal K_1, \mathcal K_2 \text{ disjoint}).
\]

\textbf{R3 (associativity / Jacobi at 4-input).}
\[
\ell_2(\ell_2(a, b), c) + (-1)^{|a|(|b| + |c|)} \ell_2(b, \ell_2(a, c)) =
\ell_2(a, \ell_2(b, c)).
\]

\textbf{R4 ($Q$-differential / BRST).}
\[
Q(\cO_{a, \beta}) = \cO_{a, \bar\partial\beta} + (\text{BRST cocycle tail}),
\]
with the cocycle tail coming from the $[\cA, \cA]$-term of
$S_{\mathrm{cl}}$.

\textbf{R5 (BV Laplacian / $\hbar$-quantum).}
\[
\Delta(\cO_{a, \beta} \otimes \cO_{b, \gamma}) \;=\;
\kappa^{ab}_\fg \, \langle \beta, \gamma \rangle_{P_{\BM}}
\cdot \mathbf{1},
\]
where
$\langle \beta, \gamma\rangle_{P_{\BM}} = \int_{\CC^3 \times \CC^3}
P_{\BM}(z, w) \beta(z) \gamma(w)$.

\textbf{R6 (shuffle-Koszul).} For $n$-fold products, the
normal-ordered shuffle rule:
\[
\cO_1 \cdot (\cO_2 \cdot \cO_3) \;=\; \sum_{\sigma \in \Sh(1,2)}
\epsilon(\sigma) (\cO_{\sigma(1)} \cdot \cO_{\sigma(2)}) \cdot
\cO_{\sigma(3)}
\]
up to $\bar\partial$-exact corrections.

\textbf{R7 (formality on $\CC^3$).}
\[
\ell_n^{E_3, \min} = 0 \text{ for } n \ge 3
\]
on flat $\CC^3$, witnessed by Kontsevich--Soibelman homotopy transfer
(Theorem~\ref{opus12:thm:ln-graph}).

\textbf{R8 (anomaly constraint).}
\[
\kappa_{\mathrm{anom}}(\CC^3, \fg) = 0,
\]
since $\chi_\top(\CC^3) = 0$ (in the Costello--Li compactly-supported
regularisation scheme); anomaly is trivial on flat $\CC^3$ for any
$\fg$.

\textbf{R9 (wave-function renormalisation).} $Z^{(1)}(\fg) = 1 -
\hbar C_2(\fg) (4\pi)^{-3} \log(L/\varepsilon)$, a cohomologically
trivial rescaling of the identity observable.

\textbf{R10 (Dualizability ${}^\natural$).} The $E_3$-trace
$\Tr_{E_3}: \Obs_{\hCS}(\CC^3) \to \CC$ via integration against
$S^5 \subset \partial \overline{\Conf}_2$ is nondegenerate on the
abelian sector $\fg = \fu(1)$; fails on non-abelian via
infinite-dim $\HH^*_{E_3}$ (Theorem~\ref{opus12:thm:dualizability}).
\end{theorem}

\begin{proof}[Proof sketch for each relation]
R1: Theorem~\ref{opus12:thm:ope}.
R2: Theorem~\ref{opus12:thm:commutativity}.
R3: Theorem~\ref{opus12:thm:mv-assoc} or
Theorem~\ref{opus12:thm:assoc-stokes}.
R4: Definition~\ref{opus12:def:bv-datum} + classical BV equation
$Q = \bar\partial + \{S_{\mathrm{cl}}, -\}_{\BV}$.
R5: Theorem~\ref{opus12:thm:ope} (equivalent form).
R6: Remark~\ref{opus12:rem:normal-ordered} + Wick's theorem.
R7: Theorem~\ref{opus12:thm:ln-graph} (Kontsevich--Soibelman on
formal CY$_3$).
R8: Proposition~\ref{opus12:prop:wheel-anomaly} at $\chi_\top = 0$.
R9: Theorem~\ref{opus12:thm:one-loop-coeff}.
R10: Theorem~\ref{opus12:thm:dualizability} below.
\end{proof}

\subsection{Deformation theory and dualizability}
\label{opus12:sec:def-dual}

\subsubsection{First-order deformations}

\begin{theorem}[First-order deformation moduli]\ClaimStatusTheorem
\label{opus12:thm:deformation-moduli}
The deformation moduli of $\Obs_{\hCS}(\CC^3)$ as an $E_3$-algebra
is
\[
\mathrm{Def}(\Obs_{\hCS}) \;=\; \HH^*_{E_3}(\Obs_{\hCS},
\Obs_{\hCS})[3].
\]
First-order deformations for simple $\fg$:
\[
T_0 \mathcal M = H^{0,3}_{\bar\partial, \mathrm{c}}(\CC^3)
\otimes \mathrm{Sym}^2(\fg^\vee)^{\mathrm{inv}}
\;\cong\; \CC \cdot \kappa_\fg,
\]
a single direction spanned by the Killing form. This matches the
Yangian deformation $Y_{\epsilon_1, \epsilon_2, \epsilon_3}(\fg)$
modulo the CY slice $\epsilon_1 + \epsilon_2 + \epsilon_3 = 0$
(wn:thm:plat-first-order-moduli).
\end{theorem}

\begin{proof}[Proof of \ref{opus12:thm:deformation-moduli}]
$\HH^*_{E_3}(\Obs_{\hCS}, \Obs_{\hCS})$ is the Hochschild cochain
complex of the $E_3$-algebra, with $[3]$-shift representing the
$E_3$-centrifugal direction. For $\Obs_{\hCS}(\CC^3) \cong
\Sym(\cE^\vee[1])[[\hbar]]$, the Hochschild cohomology at first order
reads
\[
\HH^1_{E_3}(\Obs, \Obs) = H^\bullet_{\bar\partial}(\CC^3)
\otimes \mathrm{Sym}^2(\fg^\vee)^{\mathrm{inv}}.
\]

On $\CC^3$ compactly supported:
$H^{0,3}_{\bar\partial, \mathrm{c}}(\CC^3) \cong \CC$ (by
Serre duality against $H^{0,0}(\CC^3) = \CC$), and
$\mathrm{Sym}^2(\fg^\vee)^{\mathrm{inv}} = \CC \cdot \kappa_\fg$
(the Killing form spans the invariant symmetric bilinear forms on
simple $\fg$). Hence $T_0 \mathcal M = \CC$, and the unique
first-order deformation direction is the Killing-form-weighted
action rescaling.

On the Yangian side, the $Y_{\epsilon_1, \epsilon_2, \epsilon_3}(\fg)$
family of deformed chiral algebras has three parameters
$(\epsilon_1, \epsilon_2, \epsilon_3)$ corresponding to the three
planes of the 6D hCS Omega-background. The CY condition
$\epsilon_1 + \epsilon_2 + \epsilon_3 = 0$ projects to a
$\CC^2$-family; modding out by the overall rescaling
$(\epsilon_i \to \lambda \epsilon_i)$ leaves a 1-dim parameter
space identified with $T_0 \mathcal M$ via the Costello--Gaiotto 2018
identification. $\square$
\end{proof}

\subsubsection{3-Dualizability}

\begin{theorem}[Dualizability of $\Obs_{\hCS}(\CC^3)$]\ClaimStatusTheorem
\label{opus12:thm:dualizability}
The $E_3$-trace on $\Obs_{\hCS}(\CC^3)$ is computed via
integration against the boundary 5-sphere
$S^5 \subset \partial \overline{\Conf}_2(\CC^3)$:
\[
\Tr_{E_3}(\cO) \;=\; \int_{S^5_\infty} \iota^* \cO,
\]
pulling back the observable along the collision map
$\iota: S^5 \hookrightarrow \Conf_2(\CC^3)$. The nondegeneracy
of $\Tr_{E_3}$ on the abelian sector
$\fg = \fu(1)$ gives 3-dualizability. Non-abelian 3-dualizability on
flat $\CC^3$ \emph{fails}: $\HH^*_{E_3}(\Obs_{\hCS}, \Obs_{\hCS})
= \CC[[\tau_1, \tau_2, \tau_3]]$ is infinite-dimensional
(Gwilliam--Williams 2021 Prop 5.3.2). 3-dualizability recovers on
compact CY$_3$ with finite-dim $H^\bullet(X, \omega_X)$.

$E_3$-Koszul self-duality: the operad $\mathcal D_3$ is self-dual
up to degree shift via $\mathcal D_3^! \simeq \mathrm{Lie}[2]$
(Fresse 2017 Vol I Thm 14.1.A). The strict-vs-homotopy Koszul
relation reads: Gwilliam--Williams 2021 (strict) and
Francis--Gaitsgory 2012 (homotopy) are compatible via
Fresse Thm 12.3.A + Positselski coderived/contraderived transfer.
\end{theorem}

\begin{proof}[Proof of \ref{opus12:thm:dualizability}]
Direct. The $E_3$-operad $\mathcal D_3$ has composition
$\mathcal D_3(2) \otimes \mathcal D_3(1) \to \mathcal D_3(2)$ realised
on $\Obs$ via the BV-path integration. The dualizability pairing
$\Obs \otimes \Obs \to \CC[-3]$ is the $E_3$-trace.

Explicitly, the integration against $S^5$ picks out the
\emph{residue-at-infinity} component of the factorisation product;
nondegeneracy on $\fu(1)$ follows from the finite-dim nature of the
abelian observables (a single boson per point, no non-abelian
complication). Non-abelian failure: the triple-vertex integration
over $\overline{\Conf}_3$ with generic $\fg$-structure-constants
$f^{abc}_\fg$ produces infinite-dim symmetric-algebra polynomial
tails $\CC[[\tau_1, \tau_2, \tau_3]]$, each $\tau_i$ representing a
planar loop contribution.

Recovery on compact CY$_3$: by Serre duality on $X$ compact,
$H^\bullet(X, \omega_X)$ is finite-dim, and the Hochschild cochain
complex $\HH^*_{E_3}(\Obs_X, \Obs_X)$ becomes finite-dim per
Dolbeault-cohomology degree. 3-dualizability is restored, making
$\Obs_{\hCS}(X)$ a fully-extended 3-dimensional TQFT object.
Primary: Costello--Gwilliam FA Vol 2 Ch 11; Gwilliam--Williams 2021
arXiv:2101.08385. $\square$
\end{proof}

\subsection{Cross-volume coherence: the three-volume synthesis}
\label{opus12:sec:cross-volume}

\begin{theorem}[Cross-volume coherence of
$\Obs_{\hCS}(\CC^3)$]\ClaimStatusTheorem
\label{opus12:thm:cross-vol}
The 6D hCS $E_3$-factorisation algebra $\Obs_{\hCS}(\CC^3)$
coheres across the three volumes:
\begin{itemize}
\item \textbf{Vol I scope.} $\Obs_{\hCS}(\CC^3)$ specialises under
Stage-2 factorisation $\mathrm{Sp}^{\mathrm{ch}}_{\CC, \RR}$
(factorisation homology along $\CC \subset \CC^3$, then restriction
to a reference curve) to an $E_1$-chiral algebra on the curve,
recovering the Kac--Moody vertex algebra
$V_k(\widehat{\fg})$ at level $k = \epsilon_3/\epsilon_1$ (the CY
slice $\sum \epsilon_i = 0$ reduces the three Omega-parameters to
one).
\item \textbf{Vol II scope.} The $SC^{\mathrm{ch,top}}$ bar
differential on $B(\Obs_{\hCS}(\CC^3))$ witnesses the 3D HT QFT
structure via the 6D $\to$ 5D $\to$ 4D $\to$ 3D BCOV descent. The
holomorphic factorisation $\Phi^{\mathrm{FA}}_3$ is Stage 1 of the
two-stage Vol III functor; the topological factorisation along
$\Sigma_2 = S^2$ is Stage 2.
\item \textbf{Vol III scope.} On a CY$_3$ variety $X$,
$\Phi^{\mathrm{FA}}_3(X) = \Obs_{\hCS}(X)$ is the canonical
$E_3$-holomorphic factorisation algebra. On flat $\CC^3$: the
wn:thm:plat-hCS-quantum datum with BM propagator. On $K3 \times E$:
the Künneth-multiplicative Hodge supertrace $\kappa_{\mathrm{cat}} = 0$
(AP-CY68 / AP-CY50-E24). On the quintic: nonzero
$\kappa_{\mathrm{anom}}$ requires CHSW cure.
\end{itemize}
\end{theorem}

\subsection{A first-principles-cache entry for $\Obs_{\hCS}(\CC^3)$}
\label{opus12:sec:cache}

\begin{proposition}[Cache entry: five objects, never
conflate]\ClaimStatusProposition
\label{opus12:prop:cache-five}
The programme's five-object discipline applies to 6D hCS:
\begin{enumerate}
\item \emph{Algebra} $A = \Obs^{\mathrm{cl}}_{\hCS}(\CC^3) =
\Sym(\cE^\vee[1])$ (classical).
\item \emph{Bar coalgebra} $B(A) = T^c(s^{-1} \bar A)$.
\item \emph{Dual coalgebra} $A^i = H^\star B(A)$.
\item \emph{Dual algebra} $A^! = (A^i)^\vee$ (Verdier duality).
\item \emph{Derived centre} $Z^{\mathrm{der}}_{\mathrm{ch}}(A) =
\HH^*_{E_3}(A, A)$.
\end{enumerate}
Inversion $\Omega(B(A)) = A$ is \emph{not} Koszul duality. $A^!$ is
via \emph{Verdier}. The bulk (derived centre) is via
\emph{Hochschild cochains}.

Scope-qualifying: on flat $\CC^3$, the derived centre is
$\CC[[\tau_1, \tau_2, \tau_3]]$ (infinite-dim); on compact CY$_3$,
finite-dim; the $E_3$-Koszul dual algebra $A^! \simeq \mathrm{Lie}[2]$
via Fresse 2017.
\end{proposition}

\subsection{Conclusion: the platonic object}
\label{opus12:sec:conclusion}

$\Obs_{\hCS}(\CC^3)$ is the platonic 6D hCS $E_3$-factorisation
algebra on flat complex 3-space. Its structure:

\textbf{Generators.} $\cO_{a, \beta}(\mathcal K)$ field observables
indexed by Lie-algebra index $a$, Dolbeault cochain $\beta$, compact
support $\mathcal K$.

\textbf{Propagator.} $P_{\BM}(z, w) = \frac{2}{(2\pi i)^3} \sum_k
(-1)^{k-1} \bar{(z_k - w_k)} \|z - w\|^{-6} \widehat{d\bar z_k}
\wedge dw_1 dw_2 dw_3$, the unique translation-invariant Green's
function for $\bar\partial$ on $\CC^3$ paired with
$\Omega_{\CC^3} = dw_1 dw_2 dw_3$.

\textbf{Singular OPE.} $\cA^a(z) \cA^b(w) \sim \hbar \kappa^{ab}_\fg
P_{\BM}(z, w) \cdot \mathbf{1}$ (mod $Q$-exact), derived from the
BV bracket $\{\cA(z), \cA(w)\}_{\BV} = \omega^{-1}_{\BV} =
P_{\BM}$.

\textbf{$E_3$-brackets.} $\ell_2^{E_3}(\cO_1, \cO_2) = \int
P_{\BM} \cdot \mathbf{1}$, with Koszul graded commutativity from
$\pi_1(\Conf_2(\CC^3)) = \pi_1(S^5) = 0$. Higher brackets
$\ell_n^{E_3}$ by sum-over-graphs with Koszul signs; on flat $\CC^3$
$\ell_n^{\min} = 0$ for $n \ge 3$ via Kontsevich--Soibelman
homotopy transfer (\emph{formality}).

\textbf{Associativity.} Čech--Dolbeault Mayer--Vietoris on
$\overline{\Conf}_n(\CC^3)$, equivalently Stokes on the
associahedron boundary; obstruction controlled by
$\mathrm{At}(TX)$ on compact CY$_3$.

\textbf{Commutativity.} $\pi_1(\Conf_2(\CC^3)) = \pi_1(S^5) = 0$; the
$E_3$-product is strictly graded-commutative up to
$\bar\partial$-exact corrections.

\textbf{Anomaly.} $\kappa_{\mathrm{anom}} = \hbar A(\fg)
\chi_\top(X)/(2(4\pi)^3) \|\Omega_X\|^2$; vanishes for
$\fg \in \{\mathrm{SU}(2), \mathrm{SO}(N), E_6, E_7, E_8, F_4, G_2\}$
via $d^{abc} = 0$ on these; vanishes for flat $\CC^3$ (any $\fg$)
since $\chi_\top = 0$.

\textbf{Wave-function.} $Z^{(1)}(\fg) = 1 - \hbar C_2(\fg)
(4\pi)^{-3} \log(L/\varepsilon)$; for $\mathrm{SU}(N)$, coefficient
$N/(32\pi^3)$ with the programme's conventional normalisation.
Wave-function is not anomaly (AP113 distinction).

\textbf{Dualizability.} 3-dualizable on $\fu(1)$; non-abelian
$\HH^*_{E_3}$ infinite-dim on flat $\CC^3$; 3-dualizability
recovers on compact CY$_3$.

\textbf{Deformation moduli.} $T_0 \mathcal M = H^{0,3}_c(\CC^3)
\otimes \mathrm{Sym}^2(\fg^\vee)^{\mathrm{inv}} = \CC \cdot
\kappa_\fg$, matching $Y_{\epsilon_1, \epsilon_2, \epsilon_3}(\fg)$
Yangian modulo the CY slice.

Five attack-heal-falsification cycles verify the derivation step by
step against Costello--Francis--Gwilliam 2026: propagator residue
(Cycle 1), BV-bracket OPE (Cycle 2), sum-over-shuffles
$\ell_n^{E_3}$ (Cycle 3), Čech--Dolbeault Mayer--Vietoris
associativity (Cycle 4), cubic Casimir + one-loop bubble (Cycle 5).
The generators-and-relations presentation
(Definition~\ref{opus12:def:full-generators},
Theorem~\ref{opus12:thm:relations}) can be read side-by-side with
Costello--Gwilliam FA Vol 2 Ch 5, Costello--Li 2016 §5, and
Gwilliam--Williams 2021 §5.

\subsection*{Primary literature, side-by-side}

For the reader to verify against primary sources:
\begin{itemize}
\item \textbf{Bochner--Martinelli propagator.} Range 1986,
\textit{Holomorphic Functions and Integral Representations in Several
Complex Variables}, Springer GTM 108, Ch VI §3. Classical result;
the BM formula is standard in SCV.
\item \textbf{6D hCS as a BV theory.} Costello 2013 arXiv:1303.2632;
Costello--Li 2016 arXiv:1605.09930 §2-5. The BV--BRST quantisation
of 6D hCS on $\CC^3$ and on general CY$_3$.
\item \textbf{Factorisation algebras.} Costello--Gwilliam 2017
\textit{Factorization Algebras in Quantum Field Theory} Vol 2, CUP.
Theorem 5.1 (holomorphic twist on $\CC^n$ gives $E_{2n}$
geometrically, $E_n$ after cohomology).
\item \textbf{$E_n$-Koszul self-duality.} Fresse 2017 \textit{Homotopy
of Operads and Grothendieck-Teichm\"uller Groups}, AMS Math Surv
Mono, Thm 14.1.A (strict Koszul $\mathcal D_n^! = \mathrm{Lie}[n-1]$);
Francis--Gaitsgory 2012 \textit{Selecta Math} 18 (homotopy Koszul).
\item \textbf{$L_\infty$-homotopy transfer.} Kontsevich--Soibelman
2000 \textit{Deformations of algebras over operads}, in
\textit{Conf\'erence Mosh\'e Flato}, Vol II. Explicit sum-over-trees
formula.
\item \textbf{$E_3$-trace.} Lurie HA 5.1.2.2 (Dunn additivity);
Gwilliam--Williams 2021 arXiv:2101.08385 Prop 5.3.2 (explicit
$\HH^*_{E_3}$ for Costello--Li 6D hCS).
\item \textbf{Anomaly $\kappa_{\mathrm{anom}}$.} Costello 2017
arXiv:1610.04144 §4-5; Costello 2015 arXiv:1410.5421 (Deligne
exceptional series).
\item \textbf{One-loop wave-function.} Costello 2011
\textit{Renormalization and Effective Field Theory}, AMS Math Surv
Mono, Thm 13.4.1 (compact-support existence theorem) + compactness
caveat from Costello--Gwilliam Vol 2 Prop 8.2.1 for non-compact
extension.
\item \textbf{Cubic Casimir on simple Lie algebras.} Okubo 1979
\textit{J Math Phys} 20 3253 (explicit $d^{abc}$ tensors);
Bars--Yankielowicz 1980 \textit{Phys Rev D} 21 3018 (group-theory
coefficients). Vanishing of $d^{abc}$ on exceptional and
$\mathrm{SO}$-series is standard.
\item \textbf{Yangian deformation.} Maulik--Okounkov 2012
arXiv:1211.1287 §3; Costello--Gaiotto 2018 arXiv:1810.04802 §4. The
4d-3d chiral descent matches the Omega-background on the 6D hCS
side to the Yangian deformation parameters on the curve.
\end{itemize}

\subsection*{Cross-volume references}

\begin{itemize}
\item Vol I cache entries 446, 447, 448, 449, 450 (6D hCS audit
series) and AP929-932 (Wave-14 cross-volume).
\item Vol II `bv_brst.tex` Thm `thm:6dhcs-one-loop-anomaly`, Cor
`cor:6dhcs-deligne-cancellation`, and
`thqg_perturbative_finiteness.tex` §`sec:pf-hcs-bv-feynman-trace`.
\item Vol III `chapters/theory/cy_to_chiral.tex` for the two-stage
factorisation $\Phi_d = \mathrm{Sp}^{\mathrm{ch}}_{\Sigma_{d-1}, C}
\circ \Phi^{\mathrm{FA}}_d$; `chapters/examples/cy_d_kappa_stratification.tex`
for the CY$_d$ $\kappa_\bullet$ table.
\item Cross-volume AP ledger entries AP-CY50-E1 through AP-CY50-E24
(the 2026-04-22 6D hCS audit).
\end{itemize}

% End of Agent-12 deliverable. Author: Raeez Lorgat.
